%! AP Statistics — Unit 2, Lesson 9
\documentclass[10pt]{article}
\usepackage[margin=0.6in]{geometry}
\usepackage{xcolor}
\usepackage[most]{tcolorbox}
\usepackage{enumitem}
\usepackage{multicol}
\usepackage{amsmath,amssymb}
\usepackage{array}
\usepackage{tabularx}
\tcbuselibrary{breakable}

\definecolor{sky}{HTML}{EAF3FB}
\definecolor{navy}{HTML}{1F3A5F}
\definecolor{redbox}{HTML}{FCECEC}
\definecolor{greenbox}{HTML}{EDF8F0}
\definecolor{goldbox}{HTML}{FFF7E6}
\newtcolorbox{skillbox}[2][]{
    colback=#2, colframe=navy, boxrule=0.8pt, arc=2mm,
    left=2mm,right=2mm,top=1.5mm,bottom=1.5mm,
    fonttitle=\bfseries, title={#1}, halign title=center, breakable
}
\setlist[itemize]{leftmargin=1.2em,itemsep=0.35em,topsep=0.35em}
\setlist[enumerate]{leftmargin=1.4em,itemsep=0.35em,topsep=0.35em}
\newcolumntype{C}[1]{>{\centering\arraybackslash}p{#1}}
\newcolumntype{Y}{>{\centering\arraybackslash}X}

\newcommand{\CourseName}{AP Statistics}
\newcommand{\SchoolYear}{2026--2027}
\newcommand{\MeetingLength}{55 minutes}
\newcommand{\UnitNumberName}{Unit 2: Exploring Two-Variable Data \quad}
\newcommand{\LessonNumberName}{Lesson 2.9: Analyzing Departures from Linearity}

\begin{document}
\begin{center}
    {\Large\bfseries \CourseName: \SchoolYear\ (\MeetingLength) \\
    {\small\bfseries \UnitNumberName \LessonNumberName}}
\end{center}

\begin{tcolorbox}[colback=sky,colframe=navy,boxrule=0.9pt,arc=2mm,
    left=2mm,right=2mm,top=1.5mm,bottom=1.5mm]
    \textbf{Primary Objective:} Students will identify nonlinear patterns in scatterplots and
    residual plots, and will apply logarithmic and power transformations to linearize data.
    Students will fit a linear regression to transformed data, back-transform predictions, and
    assess the quality of the transformed model (Big Idea: DAT; AP Skill 2).
\end{tcolorbox}

\begin{skillbox}[Priority Ideas \& Skills]{goldbox}
\begin{minipage}[t]{0.48\linewidth}
    \textbf{AP Skill 2 --- Data Analysis}
    \begin{itemize}
        \item Identify when a linear model is inappropriate using scatterplots and residual plots.
        \item Apply the log transformation to $y$, to $x$, or to both to achieve linearity.
        \item Fit a regression line to transformed data; write the equation in terms of the transformed variable.
        \item Back-transform to make predictions on the original scale.
    \end{itemize}
\end{minipage}
\hfill
\begin{minipage}[t]{0.48\linewidth}
    \textbf{Key Understandings}
    \begin{itemize}
        \item Transformations change the scale of the variable, not the data values themselves.
        \item After a log transformation, the regression model is exponential or power in the original variables.
        \item The residual plot for the transformed model should show random scatter --- that confirms the transformation worked.
        \item Back-transformation reverses the log: $y = 10^{\hat{y}'}$ (base-10 log) or $y = e^{\hat{y}'}$ (natural log).
    \end{itemize}
\end{minipage}
\end{skillbox}

\begin{skillbox}[{Vocabulary, Concepts, \& Theorems}]{greenbox}
\begin{minipage}[t]{0.48\linewidth}
    \begin{itemize}
        \item \textbf{Nonlinear relationship}: a pattern in a scatterplot that is better described by a curve than a line
        \item \textbf{Transformation}: a mathematical operation applied to a variable to change its distribution or relationship; common transformations: $\ln y$, $\log_{10} y$, $\sqrt{y}$, $y^2$
        \item \textbf{Exponential model}: $y = a \cdot b^x$; linearized by $\log y = \log a + (\log b)\,x$
    \end{itemize}
\end{minipage}
\hfill
\begin{minipage}[t]{0.48\linewidth}
    \begin{itemize}
        \item \textbf{Power model}: $y = a\,x^b$; linearized by $\log y = \log a + b\,\log x$
        \item \textbf{Back-transformation}: reversing the log to return to original units (e.g., $\hat{y} = 10^{\widehat{\log y}}$)
        \item \textbf{Linearization check}: construct residual plot of transformed data; random scatter confirms the transformation was successful
    \end{itemize}
\end{minipage}
\end{skillbox}

\begin{skillbox}[Activate Prior Knowledge \& Spiral Review]{sky}
\begin{minipage}[t]{0.48\linewidth}
    \textbf{Quick Review (3 min)}\\
    Display a residual plot with a clear curved pattern (from Lesson 2.7). Ask: ``What does this tell us?'' (Linear model is not appropriate.) Bridge: instead of abandoning regression, we can \textit{transform} the data to make the relationship linear, then apply the tools we already know.
\end{minipage}
\hfill
\begin{minipage}[t]{0.48\linewidth}
    \textbf{Spiral Connections}
    \begin{itemize}
        \item Lesson 2.7: nonlinear residual plots are the diagnostic that motivates transformations.
        \item Lesson 2.8: regression on transformed data uses the same formulas.
        \item Unit 9 (inference for slopes) assumes linearity --- transformation satisfies this condition.
    \end{itemize}
\end{minipage}
\end{skillbox}

\begin{skillbox}[Hook \& Lesson]{sky}
\begin{minipage}[t]{0.48\linewidth}
    \textbf{Hook (5 min)}\\
    Display the growth of a bacterial colony over time (exponential data). Ask: ``Can we fit a straight line? Look at the residual plot.'' Answer: no. Now take $\log(\text{count})$ and replot. The curved scatterplot becomes linear. Ask: ``What changed?'' --- motivates today's lesson.
\end{minipage}
\hfill
\begin{minipage}[t]{0.48\linewidth}
    \textbf{Lesson Flow (18 min)}
    \begin{enumerate}
        \item Identify the curved pattern in the original scatterplot and residual plot.
        \item Apply $\log_{10}$ to $y$; replot and note the linearization.
        \item Fit a regression to $(\log y, x)$; write the equation.
        \item Back-transform: convert the equation back to the original exponential form.
        \item Verify: inspect the residual plot of the transformed data.
    \end{enumerate}
\end{minipage}
\end{skillbox}

\begin{skillbox}[Explicit Instruction \& Active Monitoring]{sky}
\begin{minipage}[t]{0.48\linewidth}
    \textbf{Direct Instruction (10 min)}
    \begin{itemize}
        \item Exponential rule: log transform the $y$ values only; regression on $(\log y$ vs.\ $x)$ gives $\widehat{\log y} = a + bx$; back-transform: $\hat{y} = 10^{a+bx}$.
        \item Power rule: log transform both $x$ and $y$; regression on $(\log y$ vs.\ $\log x)$.
        \item AP exam emphasis: students must (1) identify nonlinearity, (2) state what transformation was used, (3) write the regression equation in transformed terms, and (4) back-transform for a prediction.
    \end{itemize}
\end{minipage}
\hfill
\begin{minipage}[t]{0.48\linewidth}
    \textbf{Active Monitoring}
    \begin{itemize}
        \item Common error: making a prediction using the transformed equation without back-transforming (e.g., reporting $\widehat{\log y} = 2.3$ instead of $\hat{y} = 10^{2.3} \approx 200$).
        \item Check: Is the residual plot for the transformed model random scatter? If not, the transformation was insufficient.
        \item Probe: ``Which variable did we transform? $x$, $y$, or both?'' --- determines which model type (exponential vs.\ power) applies.
    \end{itemize}
\end{minipage}
\end{skillbox}

\begin{skillbox}[Group Work \& Differentiation]{redbox}
\begin{minipage}[t]{0.48\linewidth}
    \textbf{Group Activity (8--10 min)}\\
    Groups receive two datasets: one that needs a log-$y$ transformation (exponential) and one that needs a log-log transformation (power). Groups:
    \begin{enumerate}
        \item Identify the nonlinearity from the scatterplot.
        \item Apply the appropriate transformation and construct the new scatterplot.
        \item Fit the regression; write the equation.
        \item Back-transform and make a prediction on the original scale.
    \end{enumerate}
\end{minipage}
\hfill
\begin{minipage}[t]{0.48\linewidth}
    \textbf{Differentiation}
    \begin{itemize}
        \item \textit{Support}: Provide a step-by-step transformation guide with the two cases clearly labeled; students follow steps with provided log values.
        \item \textit{Extension}: Students investigate whether a square-root transformation ($\sqrt{y}$) linearizes a different dataset better than a log transformation, and justify their conclusion using residual plots.
        \item \textit{Technology}: Use a graphing calculator's logarithmic and power regression options (ExpReg, PwrReg) to compare with the linearization approach.
    \end{itemize}
\end{minipage}
\end{skillbox}

\begin{skillbox}[Individual Work \& Assessment]{redbox}
\begin{minipage}[t]{0.48\linewidth}
    \textbf{Exit Ticket (5 min)}\\
    A $\log_{10}$ transformation of $y$ linearizes the data. The regression of $\log y$ on $x$ gives: $\widehat{\log y} = 1.2 + 0.15x$. Students: (1) write the back-transformed exponential equation, (2) predict $y$ when $x = 4$, (3) explain one way to verify the transformation was effective.
\end{minipage}
\hfill
\begin{minipage}[t]{0.48\linewidth}
    \textbf{AP-Style Check}\\
    \textit{``A scatterplot of [y] vs.\ [x] shows a curved pattern. After applying a log transformation to $y$, the scatterplot of $\log y$ vs.\ $x$ appears linear with $r^2 = 0.93$. Write the regression equation for $\log y$ on $x$ and use it to predict $y$ when $x = 10$.''}\\[4pt]
    Assess: correct equation in transformed form; correct back-transformation; numerical answer with context.
\end{minipage}
\end{skillbox}

\begin{skillbox}[Reinforcement \& Extension]{goldbox}
\begin{minipage}[t]{0.48\linewidth}
    \textbf{Homework / Practice}
    \begin{itemize}
        \item Full transformation analysis for two datasets: identify nonlinearity, apply transformation, fit regression, back-transform, assess with residual plot.
        \item Personal Progress Check 2 (AP Classroom) assigned for Unit 2 review.
    \end{itemize}
\end{minipage}
\hfill
\begin{minipage}[t]{0.48\linewidth}
    \textbf{Extension / Preview}
    \begin{itemize}
        \item \textit{Extension}: Research how epidemiologists use log-transformed data when modeling the spread of infectious diseases. What does the slope mean in that context?
        \item \textit{Unit 3 Preview}: Next unit shifts from analyzing data to \textit{collecting} it. The methods we use to collect data determine whether our conclusions can be generalized (from a sample to a population) or whether causation can be established (from an experiment).
    \end{itemize}
\end{minipage}
\end{skillbox}

\end{document}
